% k_paper.tex -- manuscript for Methodology (European Journal of Research
% Methods for the Behavioral and Social Sciences, PsychOpen GOLD/ZPID).
% See notes/referencias_pendientes.md and metodo_borrador.md for the
% verified Spanish-language source material this is drafted from -- see
% [[project-responsecategories-normality-study]] memory for the target-
% journal decision and its submission constraints (2026-09-26): 6,000-word
% hard limit (incl. references/tables/figures, max 6 figures), abstract
% <=150 words, mandatory double-blind anonymization, APA 7th throughout.
%
% Skeleton only at this stage: section headers and structure are in
% place, content is ported/translated from the Spanish drafts
% incrementally. Word-count triage against the 6,000-word cap happens
% once content is in, not before.
%
% Modeled directly on the proven double-blind APA7 pattern already used
% for FiabilityLab's BRM submission
% (../../FiabilityLab-CHJS/paper/paper_brm.tex): apa7.cls with
% [man,mask,donotrepeattitle], biblatex-apa, and the \blindtext{}{} macro
% for anonymizing the repository link during review.
%
% [man]: manuscript/submission mode (double-spaced, running head, single
% column), not [jou] (the final typeset format produced after
% acceptance).
% [mask]: apa7's built-in double-blind option -- suppresses author name/
% affiliation/running-head-by-author from the compiled PDF. Methodology
% additionally wants a SEPARATE title-page file at first submission (not
% just a masked single file) -- that separate title-page .tex is created
% later, once we're ready to submit; this file is the anonymized main
% text.
% [donotrepeattitle]: avoids the title repeating as a running header
% immediately above the body on page 1.
\documentclass[man,mask,donotrepeattitle]{apa7}

\usepackage{amsmath}
\usepackage[american]{babel}
\usepackage{csquotes}
\usepackage[backend=biber,style=apa,sortcites=true,sorting=nyt]{biblatex}
\DeclareLanguageMapping{american}{american-apa}
\addbibresource{k_paper.bib}

% Same blinding helper as paper_brm.tex: #1 = real content (shown once
% [mask] is removed for the final, identified version); #2 = placeholder
% shown during blind review (temporary anonymized link, per Methodology's
% author guidelines on anonymizing external materials).
\makeatletter
\newcommand{\blindtext}[2]{\@ifundefined{apaSeven@maskauthoridentity}{#1}{#2}}
\makeatother

\title{When Does the Number of Response Categories Matter? A Monte Carlo and Composite-Score Study of Its Interaction With Sample Size and Non-Normality Severity Across Eleven Normality Tests}
\shorttitle{RESPONSE CATEGORIES AND NORMALITY TEST POWER}

\authorsnames{Arqu\'imedes Chac\'on}
\authorsaffiliations{Universidad Cat\'olica Andr\'es Bello, Caracas, Venezuela}
\authornote{Correspondence concerning this article should be addressed to
  Arqu\'imedes Chac\'on, Universidad Cat\'olica Andr\'es Bello, Caracas,
  Venezuela. Email: arquimedeschacon@gmail.com}

\leftheader{Chac\'on}

% TODO: draft to <=150 words once Results/Discussion content is final --
% see notes/resultados_borrador.md and the practical-utility synthesis in
% notes/DESIGN.md Section 19 for the source material.
\abstract{[Abstract -- max. 150 words. Draft after Results/Discussion are
ported and trimmed to the 6,000-word cap; source material: the
hypothesis-by-hypothesis synthesis in notes/resultados\_borrador.md and
the practical-utility paragraph in notes/DESIGN.md Section 19.]}

\keywords{Likert scale, response categories, normality tests, Monte Carlo simulation, non-normality severity, composite score}

\begin{document}
\maketitle

\section{Introduction}
\label{sec:introduction}
% Port/translate from notes/introduccion_borrador.md (v4). Source
% structure (6 paragraphs + closing, see that file's own "Que cambio"
% notes for revision history):
%   1. Why response-category count (k) is an open, unsettled design
%      decision (Dillman et al., 2009; Foddy, 1994; Pearse, 2011;
%      Lozano, Garcia-Cueto & Muniz, 2008).
%   2. What is known about k's effect on downstream statistical
%      behavior -- fragmented, rarely aimed at normality testing
%      specifically (Weijters et al., 2010; Kurtulus et al., 2026;
%      Ho & Yu, 2014; Sathyanarayana, 2026; Cain et al., 2017).
%   3. The parallel tradition of comparing normality-test power, which
%      varies distribution shape but never k as an independent factor
%      (Razali & Wah, 2011; Shapiro & Wilk, 1965; Yap & Sim, 2011;
%      Bantu et al., 2025; Matamoros-Mota, 2026; Kamath et al., 2025).
%   4. The specific, uncovered gap: how k, n, and non-normality severity
%      jointly affect normality-test power.
%   5. The two-component study design (simulated niveles block + 6 real
%      instruments), stated explicitly.
%   6. [Closing with explicit study objectives -- still a placeholder in
%      the Spanish draft too, pending final Results section order.]
%
% NOTE: several citations here (Pearse 2011, Lozano et al. 2008,
% Weijters et al. 2010, Kurtulus et al. 2026, Ho & Yu 2014,
% Sathyanarayana 2026, Razali & Wah 2011, Yap & Sim 2011, Bantu et al.
% 2025, Matamoros-Mota 2026, Kamath et al. 2025) are NOT YET in
% k_paper.bib -- verify and add before citing with \parencite{}.

\section{Method}
\label{sec:method}

\subsection{General Design}
\label{sec:design}
% Port from notes/metodo_borrador.md Section 2.1.

\subsection{Real Psychometric Instruments}
\label{sec:instruments}
% Port Table 1 (instrument characteristics: k, items, dimensions, N,
% alpha, reverse-scored items, dataset source) and surrounding prose from
% notes/metodo_borrador.md Section 2.2. Dataset citations already
% verified and in k_paper.bib: \parencite{osp2014rse, osp2019mach,
% yamada2021, osp2014hexaco, davis2021, osp2016rwas}. Instrument-origin
% citations also ready: \parencite{rosenberg1965, christiegeis1970,
% cutronarussell1987, ashtonleegoldberg2007, snyder1991, altemeyer1981}.
% SPS-10 dimensionality citation: \parencite{sischka2025}. RSE alpha
% support: \parencite{blascovichtomaka1993, rosenberg1986}.

\subsection{Niveles Block: Simulating Non-Normality Severity}
\label{sec:niveles}
% Port from notes/metodo_borrador.md Section 2.3, including Table 2
% (five severity levels anchored to Cain et al. 2017 percentiles) and
% the calibration procedure (free-lambda Nelder-Mead, 16 starts,
% N=150,000 common-random-number replicates). Cite \parencite{cain2017}.

\subsection{Real-Data Block: \emph{m}-out-of-\emph{N} Subsampling}
\label{sec:realdata}
% Port from notes/metodo_borrador.md Section 2.4. Cite
% \parencite{politisromanowolf1999}.

\subsection{Normality Tests Evaluated}
\label{sec:tests}
% Port from notes/metodo_borrador.md Section 2.5 (eleven-test battery).
% Citations ready: \parencite{shapirowilk1965, andersondarling1954,
% lilliefors1967, jarquebera1980, dagostinopearson1973, csorgofaraway1996,
% shapirofrancia1972, pearson1900, anscombeglynn1983, anaratschwender2026}.
% Epps-Pulley has no formal citation in the Spanish draft either --
% resolve before finalizing this subsection.

\subsection{Implementation and Reproducibility}
\label{sec:implementation}
% Port from notes/metodo_borrador.md Section 2.6.

\subsection{Data and Code Availability}
\label{sec:availability}
% Port from notes/metodo_borrador.md Section 2.7. Use \blindtext{}{} for
% the repository link, same pattern as the Software/Data/Code
% Availability sections below.

\section{Results}
\label{sec:results}
% Port from notes/resultados_borrador.md: Tables 3-7, Figure 1, the H1-H5
% hypothesis tests, the validation against real data (r = .981), and the
% eleven-test comparison. Methodology allows a maximum of six figures --
% check notes/resultados_borrador.md's figure count against that cap
% before porting.

\section{Discussion}
\label{sec:discussion}
% Port/expand from the practical-utility synthesis in
% notes/DESIGN.md Section 19, plus limitations (e.g., the
% percentile-90-vs-95 degenerate-calibration lesson from Section 17,
% worth stating explicitly as a methodological contribution in its own
% right).

\section*{Data Availability}
% TODO: finalize once the repository's public/anonymized status is
% decided (see the "Pendiente" note in notes/metodo_borrador.md Section
% 2.7's surrounding notes). Placeholder pattern below matches
% FiabilityLab's Anonymous GitHub approach.
The full per-cell simulation and subsampling results analyzed in this
study are available at \blindtext{\url{https://github.com/ArchieJamDev/ResponseCategories-Normality-Study}}{an anonymized mirror of the repository (identifying information removed for double-blind review; the original link will be restored upon acceptance)}. Raw data for the six real instruments are publicly available from their original sources, cited where first used in Section~\ref{sec:instruments}.

\section*{Code Availability}
The complete code (real-data extraction, niveles-block calibration and
simulation, real-data-block subsampling, results aggregation, and the
GitHub Actions configuration that runs every step) is available at
\blindtext{\url{https://github.com/ArchieJamDev/ResponseCategories-Normality-Study}}{an anonymized mirror of the repository (identifying information removed for double-blind review; the original link will be restored upon acceptance)}.

\section*{Use of Generative Artificial Intelligence}
Generative artificial intelligence (Claude, Anthropic) was used to
implement and run R code -- including the simulation and real-data
analysis scripts reported in this article -- and as a writing assistant
for the text, always under the author's direct review and direction.
Conceptual and methodological decisions -- what to simulate, how to
design each block, which real instruments to use, how to interpret each
finding, and the article's overall structure and argument -- were the
author's.

\section*{Declarations}

\textbf{Funding.} The author received no specific funding for this work.

\textbf{Competing interests.} The author declares no competing interests.

\textbf{Ethics approval.} Not applicable. No new data were collected from
human participants; all real-data examples reuse previously published,
publicly available, de-identified datasets (see Section~\ref{sec:instruments}
and the Data Availability statement above).

\textbf{Consent to participate and consent for publication.} Not
applicable.

\textbf{Author contributions.} Single-author manuscript; A.C. conceived
the study, designed and ran the analyses, and wrote the manuscript.

\section*{Open Practices Statement}
No part of this study was preregistered; it is a methodological and
simulation study, not a confirmatory hypothesis test. All simulation and
subsampling code, and the full per-cell results, are available as
described in the Data and Code Availability statements above.

\printbibliography

\end{document}
